\section{Fazit}

Diese Masterarbeit hat erfolgreich eine zweistufige Regelungsarchitektur für ein selbstbalancierendes Fahrrad entwickelt und in der 
Webots-Simulationsumgebung validiert. Die systematische Integration einer hochfrequenten PID-Balanceregelung (500\,Hz) mit einer
 visionbasierten MPC-Spurführung (20\,Hz) stellt einen wichtigen methodischen Fortschritt in der Entwicklung autonomer Fahrradsysteme dar.

\subsection{Zentrale Erkenntnisse}

Die entwickelte duale Controller-Architektur demonstriert erstmals die praktische Machbarkeit der Kopplung von 
Balance- und Fahrtrichtungsregelung in einer integrierten Simulationsumgebung. Die erfolgreiche Portierung 
bewährter Hardware-Implementierungen von Zander (2023) und Giray (2024) zeigt die Übertragbarkeit etablierter 
Regelungskonzepte zwischen verschiedenen Plattformen auf.

Ein überraschendes und bedeutsames Ergebnis ist die durchgängige Überlegenheit der ROI-basierten Fallback-Erkennung 
gegenüber der YOLO-Segmentierung. In allen Testszenarien erwies sich die klassische Bildverarbeitung als überlegen: 
Bei der Kurvenfahrt erreichte ROI-Fallback 43,62\,s stabile Fahrt gegenüber 35,96\,s bei YOLO-PID, in der S-Kurve 
ergab sich eine 16,8\% Verbesserung der Fahrtdauer mit halbierten Querfehlern. Diese Ergebnisse zeigen, dass für 
strukturierte Simulationsumgebungen klassische Bildverarbeitung durchaus konkurrenzfähig zu modernen 
Deep-Learning-Ansätzen sein kann.

Die systematische Evaluierung verschiedener Störungsszenarien belegt die Anpassungsfähigkeit der 
Regelungsarchitektur. Das System bewältigt Seitenwind bis 1,0\,m/s und Höhenunterschiede bis 4\,m 
durch adaptive PID-Parametrierung, wobei sich Senken als weniger kritisch als Hügel erweisen.

\subsection{Methodische Beiträge}

Die Arbeit etabliert eine systematische Methodik für die risikofreie Entwicklung und Validierung 
komplexer Regelungsstrategien in Webots. Das entwickelte JSON-basierte Konfigurationssystem mit 
grafischem Interface unterstützt iteratives Parameter-Tuning, während die hochfrequente Datenerfassung 
(2\,ms Abtastung) detaillierte Regeldynamik-Analysen ermöglicht.

\subsection{Limitierungen und kritische Reflexion}

Die identifizierten Grenzen ordnen die Ergebnisse in einen realistischen wissenschaftlichen Kontext ein. 
Besonders kritisch für eine praktische Umsetzung erweisen sich die sensorischen Idealisierungen (perfekte IMU- 
und Kameradaten ohne Rauschen oder Drift), die vereinfachte ODE-basierte Physikmodellierung und die 
Echtzeitfähigkeits-Herausforderungen bei der Vision-Pipeline auf eingebetteten Systemen.

Die Hardware-Transfer-Problematik manifestiert sich auf mehreren Ebenen: von der numerischen Stabilität der 
Physiksimulation über Kommunikationslatenzen bis hin zu mechanischen Toleranzen und elektronischen Nichtlinearitäten, 
die in der Simulation vollständig abstrahiert werden.

\subsection{Wissenschaftlicher Beitrag}

Diese Arbeit leistet vier wesentliche Beiträge zur Forschung: (1) die erstmalige erfolgreiche Integration von 
Balance- und Spurführungsregelung, (2) systematische Methodenvergleiche zwischen verschiedenen Ansätzen, (3) 
umfassende Robustheitsvalidierung unter verschiedenen Störungsszenarien und (4) die Entwicklung einer modularen, 
erweiterbaren Regelungsarchitektur.

\subsection{Ausblick}

Methodisch unterstreichen die aufgezeigten Grenzen die Notwendigkeit einer mehrstufigen Validierungsstrategie, 
die von der Simulation über Hardware-in-the-Loop-Tests bis hin zu systematischen Feldversuchen reicht. Zukünftige 
Forschungsrichtungen sollten sich auf die systematische Überbrückung der Sim-to-Real-Lücke konzentrieren, 
insbesondere durch Integration hardware-spezifischer Charakteristika, erweiterte Physikmodellierung und 
Multi-Environment-Training.

Trotz der identifizierten Limitierungen bildet die vorliegende Arbeit eine fundierte Grundlage für die 
Weiterentwicklung autonomer selbstbalancierender Fahrradsysteme. Die erfolgreiche Demonstration der 
kombinierten Regelungsarchitektur und die detaillierte Dokumentation der Systemgrenzen schaffen eine 
solide Basis für zukünftige Forschungsarbeiten und unterstreichen die Bedeutung einer kritischen Reflexion 
der Übertragbarkeit von Simulationsergebnissen auf reale Anwendungen.